
\vsize 18cm
\hsize 15.4cm
\magnification=1200
\item{}
\vskip 1cm
\font\r=cmr10 scaled\magstep2

\centerline{\r INVARIANTS AND RELATED LIAPUNOV} 
\medskip
\centerline{\r FUNCTIONS FOR DIFFERENCE EQUATIONS II}

\bigskip 
\centerline{by}
\bigskip
\centerline{M.R.S.Kulenovi\'c}

\centerline{Department of Mathematics}

\centerline{University of Rhode Island}

\centerline{Kingston, RI 02881}

\vskip 1.5cm
\centerline{\bf Abstract}
\vskip 1.5cm
Consider the difference equation 
$${\bf x_{n+1}}={\bf f( x_n)},$$
where ${\bf x_n}$ is in $R^k$  and ${\bf f}: D\to D$, is continuous  
where $D \subset R^k$. 
Suppose that $I:R^k\to R$ is continuous invariant that is, $I({\bf f(x)})=I({\bf x})$ 
for every ${\bf x}\in D$. Using our previous result which states 
that the  Liapunov function is equal to  $\pm(I({\bf x}) - I({\bf p}))$, 
where ${\bf p}$ is an isolated  equilibrium (fixed) point  of this system in 
which $I$ attains the minimum or maximum value, see [K], we prove the stability of some 
difference equations that appear in mathematical physics such as: 
$$(x_{n+1} + x_{n-1})(1+{1\over 2}x_n^2) - (a+2)x_n=0,$$ 
where $a>0$ is a parameter. Using our result we obtain the Liapunov function 
$$V(x,y)=x^2y^2-2(a+2)xy+2(x^2+y^2)+a^2.$$
We also give the method for constructing the invariants for a class of 
equations that involves rational functions. 


\vskip 2cm

AMS 1980 Subject Classification: Primary-39A11, Secondary-39Al0


\vfill\eject
\vglue 2cm

{\r 1. Introduction}
\bigskip
Consider the soliton equation 
$$(x_{n+1} + x_{n-1})(1+{1\over 2}x_n^2) - (a+2)x_n=0,$$ 
where $a>0$ is a parameter, and $x_{-1}>0, x_0>0$. Recently, this equation 
attracted a lot of attention and its invariant (first integral) has been found,  
see [QRT],[RQ] and [HBQC]:
$$I(x_n,x_{n-1})= x_n^2 x_{n-1}^2 - 2(a+2)x_n x_{n-1} + 2(x_n^2 + x_{n-1}^2).$$ 
Here,  we will construct the Liapunov function for this equation using this  
invariant. 
In fact, we will show that the Liapunov function $V(x,y)$ is given by
$$V(x,y)=I(x,y)-I(p,p)=I(x,y)+a^2,$$
where 
$$p= \sqrt a,$$
is an equilibrium (fixed point) of this equation. To do this, we will use the 
following result obtained in [K]:
\bigskip
{\bf THEOREM A.} {\it Consider the  difference equation 
$${\bf x_{n+1}}={\bf f( x_n)}$$
where ${\bf x_n}$ is in $R^k$, and ${\bf f}: D\to D$ is continuous, where $D \subset R^k$. 
Suppose that $I:R^k\to R$ is an invariant that is, $I({\bf f(x)})=I({\bf x})$ 
for every ${\bf x}\in D$. If $I$ attains an isolated local minimum or maximum value 
at the  equilibrium (fixed) point ${\bf p}$  of this system, then there exists a 
Liapunov function equal to  $\pm(I({\bf x}) - I({\bf p}))$, and so the equilibrium 
${\bf p}$ is stable.} 
\bigskip
The stability considered here is taken in the sense of [KP]: 
\medskip
{\bf DEFINITION} ([KP], p.171) Let {\bf u} belong to ${\bf R}^n$ and $r>0$. The open
ball centered at {\bf u} with radius $r$ is the set
$$B({\bf u},r) = \{{\bf v} \in {\bf R}^n: |{\bf v} - {\bf u}|<r\}.$$
Let {\bf p} be a fixed point of {\bf f}. Then {\bf p} is stable provided that,
given any ball B({\bf p},$\epsilon$), there is a ball B({\bf p},$\delta$) so that
if {\bf u} is in B({\bf p},$\epsilon$), then ${\bf f}^k ({\bf p })$ is in B({\bf p},$\epsilon$)
for $k=1,2,...$. Here ${\bf f}^k ({\bf p })$ denotes $k$-th iterate of ${\bf p}$.
\medskip
{\bf REMARK 1.} This theorem shows that whenever some conditions are satisfied, 
there exists a Liapunov function which is a special invariant of considered equation. 
This emphasizes the importance of developing the methods for finding the invariants. 
Some methods has been developed in [GKL], [GM] and [QRT].  
In this paper, we will elaborate on the method developed in [QRT]. 
As is well known, see [GM] and [M], the invariant is not uniquely determined but 
rather if one invariant is given then any appropriate function (continuous, 
analytic) of this invariant is invariant itself. Furthermore, the difference equation can 
have several "independent" invariants (eg. two invariants neither of which can
be derived from the other through homeomorphic changes of variable), see [BB]. As we will 
mention later we have developed {\it Mathematica} package that can almost automatically find the
rational invariant for a given equation (if there is one) and using our result here construct the 
corresponding Liapunov function, if such function exists.
\bigskip
{\r 2. Main Result. }
\bigskip
Our main result is based on the following lemmas:
\bigskip
{\bf LEMMA 1.} {\it Consider the vectors ${\bf x}=((g(x))^2,g(x),1)^T$,  and ${\bf z}=
(z^2,z,1)^T$, and the vector-valued function 
$${\bf f(y)} = (f_1 (y), f_2 (y), f_3 (y))^T,$$
where $g,f_1,f_2,f_3$ are the given functions and 
$$z\neq g(x), \quad f_2(y) - g(x)f_3(y)\neq 0.\leqno(1)$$
Then, the mixed product of ${\bf x,z}$ and ${\bf f(y)}$ is zero, that is 
$${\bf x \times f(y) \cdot z}=0,$$ 
if and only if 
$$z={{f_1(y) - g(x)f_2(y)}\over {f_2(y) - g(x)f_3(y)}}. \leqno(2)$$}
\medskip
{\bf Proof.}  
$$0={\bf x \times f(y)\cdot z} = f_1(y)(z-g(x)) + f_2(y)((g(x))^2 - z^2) + 
f_3(y) (g(x)z^2 - z(g(x))^2)$$
$$=(z-g(x))[f_1(y) - f_2(y) (g(x) + z) + f_3(y) g(x)z].$$
If $(1)$ holds, we conclude that the last equality is equivalent to $(2)$, 
which completes the proof of lemma.
\bigskip
{\bf REMARK 2.} Taking in $(2)$ $z=x_{n+1},\quad y=x_n$, and $x=x_{n-1}$ we get 
the following difference equation 
$$x_{n+1}={{f_1(x_n) - g(x_{n-1})f_2(x_n)}\over {f_2(x_n) - g(x_{n-1})f_3(x_n)}}, \leqno(3)$$
\medskip
\noindent
while if we take $z=x_{n+1},\quad y=x_{n-1}$, and $x=x_n$ we obtain the similar 
difference equation:
$$x_{n+1}={{f_1(x_{n-1}) - g(x_{n})f_2(x_{n-1})}\over {f_2(x_{n-1}) - g(x_{n})f_3(x_{n-1})}}. \leqno(4)$$
\medskip
The next result establishes the necessary connection between the difference 
equations of the types $(3)$, and $(4)$ and their invariants.
\bigskip
{\bf LEMMA 2.} {\it Suppose that in Lemma 1 ${\bf x\neq 0,\quad f(y) = A_1 y \times A_2 y}$ where 
$${\bf y}= (f(y)^2, f(y), 1)^T,$$ 
where $f$ is a given function, and ${\bf A_i},\; i=1,2$ are given matrices 

$${\bf A}_i=\left(\matrix{a_i&b_i&c_i\cr
                  d_i&e_i&f_i\cr
                  g_i&h_i&k_i\cr}\right),$$
$i=1,2$. Then, for every vector ${\bf x}$, such that ${\bf (A_2 y)\cdot x}\neq 0$,  
the expression 
$${{{\bf (A_1 y)\cdot x}}\over {{\bf (A_2 y)\cdot x}}}, \leqno(5)$$ 
is independent of ${\bf x}$, that is:
$${{{\bf (A_1 y)\cdot x}}\over {{\bf (A_2 y)\cdot x}}} ={{{\bf (A_1 y)\cdot z}}\over {{\bf (A_2 y)\cdot z}}},\leqno(6)$$ 
for every ${\bf z}$, for which ${\bf (A_2 y)\cdot z}\neq 0$.}
\medskip
{\bf Proof.}  Using the fact that ${\bf (a\times b)\times c} = 
{\bf (a\cdot c) b - (a\cdot b)c}$, we get 
$$0={\bf x \times f(y)\cdot z} = {\bf x \cdot (f(y)\times z)} =  
{\bf x\cdot [(A_1y\times A_2y)\times z]} $$ 
$$= -{\bf x\cdot [(z\cdot A_2y) A_1y - (z\cdot A_1y) A_2y}],$$
which, by ${\bf x\neq 0}$, is equivalent to 
$${\bf (z\cdot A_2y) (A_1y\cdot x)= (z\cdot A_1y) (A_2y \cdot x)},$$
which in view of the conditions of lemma implies the condition $(6)$.
\bigskip
{\bf REMARK 3.} The expression $(5)$ is exactly an invariant and $(6)$ is the 
corresponding condition for the invariance of this expression. In fact, 
taking in $(6)$, $ y=x_n,\quad \hbox{and} \quad x=x_{n-1}$ we get a polynimal-like 
invariant:
$$I(x_n,x_{n-1})= \leqno(7)$$
$$(a_1 + Ka_2)(g(x_{n-1}))^2 (f(x_n))^2 + (b_1 + Kb_2) (g(x_{n-1}))^2 (f(x_n)) + 
(d_1 + Kd_2)(g(x_{n-1})) (f(x_n))^2 $$
$$+ (c_1 + Kc_2)(g(x_{n-1}))^2 + (e_1 + Ke_2)g(x_{n-1}) f(x_n) + (g_1 + Kg_2) (f(x_n))^2  $$
$$+ (f_1 + Kf_2) g(x_{n-1}) + (h_1 + Kh_2) f(x_n) + k_1 + Kk_2,$$
where $K$ is a constant with the value determined  by the initial points $x_0$,  
and $x_1$ as 
$$K= - {{{\bf (A_1 y_0)\cdot x_0}}\over {{\bf (A_2 y_0)\cdot x_0}}},$$
where ${\bf y_0}=((f(x_1))^2,f(x_1),1)$, and ${\bf x_0}=((g(x_0))^2,g(x_0),1)$. 
Equating the right hand side of $(7)$ to $0$ and solving it for $K$, we get the rational-like invariant. The expression $(7)$ is exactly the invariant of equation $(3)$. 

Taking in $(6)$, $ y=x_{n-1},\quad \hbox{and} \quad x=x_n$, we get a polynimal-like 
invariant:
$$I(x_n,x_{n-1})= \leqno(8)$$
$$(a_1 + Ka_2)(f(x_{n-1}))^2 (g(x_n))^2 + (b_1 + Kb_2) (g(x_n))^2 (f(x_{n-1})) + 
(d_1 + Kd_2)(g(x_n)) (f(x_{n-1}))^2 $$
$$+ (c_1 + Kc_2)(g(x_n))^2 + (e_1 + Ke_2)g(x_n) f(x_{n-1}) + (g_1 + Kg_2) (f(x_{n-1}))^2  $$
$$+ (f_1 + Kf_2) g(x_n) + (h_1 + Kh_2) f(x_{n-1}) + k_1 + Kk_2,$$
where $K$ is a constant with the value determined  by the initial points $x_0$,  
and $x_1$ in a similar way as above. The expression $(8)$ is an invariant of 
the equation $(4)$.

In the special case,  when $g(x)=x$, and $f(y)=y$, we get the rational-like equations 
$$x_{n+1}={{f_1(x_n) - x_{n-1}f_2(x_n)}\over {f_2(x_n) - x_{n-1}f_3(x_n)}}, \leqno(3')$$
and
$$x_{n+1}={{f_1(x_{n-1}) - x_{n}f_2(x_{n-1})}\over {f_2(x_{n-1}) - x_{n}f_3(x_{n-1})}}, \leqno(4')$$
with corresponding rational invariants:
$$(a_1 + Ka_2)x_{n-1}^2 x_n^2 + (b_1+Kb_2)x_n x_{n-1}^2 + (d_1+Kd_2)x_n^2 x_{n-1} + 
(c_1+Kc_2)x_{n-1}^2 \leqno(7')$$
$$+(e_1+Ke_2)x_{n-1} x_n + (g_1+Kg_2)x_n^2 + (f_1+Kf_2)x_{n-1} + (h_1+Kh_2)x_n+k_1+Kk_2,$$
and
$$(a_1 + Ka_2)x_{n-1}^2 {x_n}^2 + (b_1+Kb_2)x_n^2 x_{n-1} + (d_1+Kd_2)x_n x_{n-1}^2 + 
(c_1+Kc_2)x_n^2 \leqno(8')$$
$$+(e_1+Ke_2)x_{n-1} x_n + (g_1+Kg_2)x_{n-1}^2 + (f_1+Kf_2)x_n + (h_1+Kh_2)x_{n-1}+k_1+Kk_2,$$
respectively. This special case has been obtained in [QRT], and used in many papers 
refered in [RQ] and [HBQC]. Further simplifications can be obtained by assuming that 
the matrices ${\bf A_i}, i=1,2$ are symmetric, that is ${\bf A_i^T=A_i} , i=1,2$ 
which reduces the invariants $(7')$ and $(8')$ to the same form
$$(a_1 + Ka_2)x_{n-1}^2 x_n^2 + (b_1+Kb_2)(x_n^2 x_{n-1} + x_n x_{n-1}^2) + 
(c_1+Kc_2)(x_n^2+x_{n-1}^2) \leqno(9)$$
$$+(e_1+Ke_2)x_{n-1} x_n  + (f_1+Kf_2)(x_n + x_{n-1}) + k_1+Kk_2.$$
Some of these results can be formulated as a theorem.
\bigskip        
{\bf THEOREM.} {\it The expression $(7)$ (resp. $(7')$ is an invariant of the difference 
equation $(3)$ (resp. $(3')$). Similarly, the expression $(8)$ (resp. $(8')$ is an 
invariant of the difference equation $(4)$ (resp. $(4')$).}
\bigskip
{\r 3. Applications.}
\bigskip
In this part, we will use our result to find the invariants of several well known 
equations such as Lyness' equation, May's host parasitoid model, soliton model, 
Gumovski-Mira's equation as well as some new types of equations. We will also 
combine some of these results with Theorem A, to obtain the Liapunov function 
and prove the stability of their equilibria.
\bigskip
{\bf Example 1.} (see [GKL]):
\medskip
$$x_{n+1}={{a+b x_n + cx_n^2}\over {(c + dx_n+ ex_n^2)x_{n-1}}},\leqno(10)$$  
where $a,b,c,d,e$ are positive parameters. This equation can be obtained by 
taking $a_1=b_1=c_1=f_1=k_1=0$, and $a_2=e, b_2=d, c_2=c, f_2=b, k_2=a$ in the 
equation $(3')$, and the corresponding invariant $(7')$ which results in an 
invariant of the form (it is solved for $K$):

$$I(x_n,x_{n-1})={a\over {x_n x_{n-1}}} + b({1\over x_n} + {1\over x_{n-1}}) + 
c({x_{n-1}\over x_n} + {x_n\over x_{n-1}}) + d(x_n + x_{n-1}) + ex_n x_{n-1}.\leqno(11)$$

Equilibrium $p$ of this equation satisfies $ep^4+dp^3-bp-a=0.$ Now, we try to 
apply Theorem A. The necessary conditions for the extremum give:
$${\partial I \over \partial x}\equiv -{{a}\over {x^2y}} - {b\over {x^2}} + 
c({1\over y} - {y\over{x^2}}) + d + ey =0,$$
and
$${\partial I \over \partial x}\equiv -{{a}\over {xy^2}} - {b\over {y^2}} + 
c({1\over x} - {x\over{y^2}}) + d + ex =0,$$
which after some manipulations leads to $x=y$, and to the equation 
$$ex^4+dx^3-bx-a= 0.$$ 
This shows that the 
critical points are exactly equilibrium points. Now, we will check the Hessian 
at the positive critical point $(p_+,p_+)$, which will be denoted simply as $(p,p)$: 
$$H=\left(\matrix{A&B\cr
                  B&A\cr}\right),$$
where 
$$A={\partial^2 I\over \partial x^2}(p,p)={2\over p^4}(a+bp+cp^2)>0,$$
and
$$B={\partial^2 I\over {\partial x \partial y}}(p,p)={{ep^4-2cp^2+a}\over p^4}. $$
Now,  
$$det\; H=A^2 - B^2 = (A-B)(A+B)={1\over p^7}(dp^2+4cp+b)(ep^4+2bp+3a)>0,$$ 
which shows that the Hessian $H$ is positive definite at the point $(p,p)$,  
which means that the invariant $I$ attains its minimum at $(p,p)$. Thus,   
$$\min \{I(x,y): (x,y)\in D\}=I(p,p),$$ 
and so by our theorem $V(x,y)=I(x,y)-I(p,p)$ which shows  that $p$ is stable. 
\noindent
Liapunov function is given by: 
$$V(x,y)=I(x,y)-I(p,p)= I(x,y)-{{dp^3+2cp^2+3bp+2a}\over {p^2}}.$$ 
In the special case, where $a_2=0, b_2=c_2=1, k_2=a, f_2=a+1$ we obtain Lyness' 
equation 
$$x_{n+1}={{a + x_n}\over x_{n-1}},$$  
with invariant 
$$I(x_n,x_{n-1})=(1+{1\over x_n})(1+{1\over x_{n-1}})(a+x_n+x_{n-1}),$$
whose stability was established in [K]. 

In the special case where $a_2=f_2=k_2=0, b_2=c_2=1$, we get a special case 
of May's host-parazitoid model 
$$x_{n+1}={{x_n^2}\over {(1+x_n)x_{n-1}}}$$
which invariant is 
$$I(x_n,x_{n-1})=x_n + x_{n-1} + {{x_n}\over {x_{n-1}}} + {{x_{n-1}}\over {x_n}}.$$

In the special case $c=d=0$, and $a=-d, b=d^2, e=1$, we get well known Yang-Baxter's 
equation, see [QRT]
$$x_{n+1}={{d^2 x_n - d}\over {x_{n-1}x_n^2}},$$  
which invariant is given by 
$$I(x_n,x_{n-1})\equiv x_n^2x_{n-1}^2 + Kx_nx_{n-1} + d^2(x_n+x_{n-1})-d=0.$$

This equation is a special case of the class of equations that can be obtained 
from $(2)$ by taking that $f_2(y)\equiv 0$, that is the equations of the type 
$$x_{n+1}=-{{f_1(x_n)}\over {f_3(x_n) x_{n-1}}}. $$
Assuming that $f_1(y)$ and $f_3(y)$ are polynomials of the above type, we will  
get again equation $(10)$.

\bigskip
{\bf Example 2.} Gumovski-Mira equation (see [GM] and [La]):
\medskip
$$x_{n+1}={{2x_n}\over {A^2-x_n^2}} - x_{n-1},$$  
where $a$ is a positive parameter. This equation is a special case of the class 
of equations that can be obtained from $(2)$ by taking that $f_3(y)\equiv 0$, that 
is the equations of the type 
$$x_{n+1}={{f_1(x_n)}\over {f_2(x_n)}} - x_{n-1}. \leqno(12)$$
Assuming that $f_1(y)$ and $f_2(y)$ are polynomials of the above type, and taking 
that $a_1=b_1=c_1=e_1=f_1=0$  and droping the indices in other coefficients, we 
get the following equation 
$$x_{n+1}={{bx_n^2+ex_n+f}\over {ax_n^2+bx_n+c}} - x_{n-1}, \leqno(13)$$
with an invariant:
$$ Ka_2 x_{n-1}^2 x_n^2 + Kb_2 x_n x_{n-1}^2 + (d_1+Kd_2)x_n^2 x_{n-1} + 
Kc_2 x_{n-1}^2 $$
$$+Ke_2 x_{n-1} x_n + (g_1+Kg_2)x_n^2 + Kf_2 x_{n-1} + (h_1+Kh_2)x_n+k_1+Kk_2=0,$$
or in symmetric form $(9)$
$$ Ka_2 x_{n-1}^2 x_n^2 + Kb_2 (x_n^2 x_{n-1} + x_n x_{n-1}^2) + 
Kc_2 (x_n^2+x_{n-1}^2) $$
$$+Ke_2 x_{n-1} x_n  + Kf_2 (x_n + x_{n-1}) + k_1+Kk_2=0, $$
or
$$ a_2 x_{n-1}^2 x_n^2 + b_2 (x_n^2 x_{n-1} + x_n x_{n-1}^2) + 
c_2 (x_n^2+x_{n-1}^2) +e_2 x_{n-1} x_n  + f_2 (x_n + x_{n-1}) = \widetilde K, \leqno(14)$$
where $\widetilde K$ is the constant determined by the initial values $x_0$, and $x_1$. 

Taking in $(13)$, and $(14)$  $b=f=0$, and $a=-1, c=A^2, a=1$ we get Gumovski-Mira 
equation, with an invariant of the form:
$$I(x_n,x_{n-1})=x_n^2 x_{n-1}^2 - A^2(x_n^2 + x_{n-1}^2) + 2x_nx_{n-1}.$$

\bigskip
{\bf Example 3.} (see [QRT]):
\medskip

Taking in $(13)$  and $(14)$ $b=f=0$ and $e=\omega +2, a=1/2, c=1$ we get 
the soliton's equilibrium equation from [QRT] 
$$x_{n+1}={{(\omega +2)x_n}\over {1+{x_n^2}/ 2}} - x_{n-1},$$  
with an invariant of the form:
$$I(x_n,x_{n-1})=x_n^2 x_{n-1}^2 + 2(x_n^2 + x_{n-1}^2) - 2(\omega +2)x_nx_{n-1}.$$
Equilibrium $p$ of this equation satisfies $p^2-\omega=0$, that is $p=\pm \sqrt{\omega}$,  
if we additionaly assume that $\omega\geq 0$. Now, we try to 
apply Theorem A. The necessary conditions for the extremum give:
$${\partial I \over \partial x}\equiv 2xy^2-2(\omega +2)y + 4x =0,$$
and
$${\partial I \over \partial x}\equiv 2x^2y-2(\omega +2)x + 4y =0,$$
which leads to $x=y$, and to the equations $x^2-\omega=0$, and $x=0$. This shows 
that the critical points are exactly the equilibrium points. Now, we will check the 
Hessian at the positive critical point $p= \sqrt{\omega}$ :

$$H=\left(\matrix{A&B\cr
                  B&A\cr}\right),$$
where 
$$A={\partial^2 I\over \partial x^2}(p,p)=2 \omega +2 >0, $$
and

$$B={\partial^2 I\over {\partial x \partial y}}(p,p)=4\omega - 2(\omega +2)=2\omega -4.$$
Now,  
$$det\; H=A^2 - B^2 = 32\omega>0,$$ 
which shows that the invariant $I$ attains its minimum at $(p,p)$. Thus,   
$$\min \{I(x,y): (x,y)\in D\}=I(p,p)=-\omega^2,$$ 
and so by our theorem $V(x,y)=I(x,y)+ \omega^2$, which shows that $p$ is stable. 
\bigskip
{\bf Example 4.} 
\medskip
Taking in $(3')$  and $(7')$ $f_1(x)=x+c-c^2, \, f_2(x)=-cx$ and $f_3(x)=-x^2$ 
we get the following equation: 
$$x_{n+1}={{cx_nx_{n-1}+x_n+c-c^2}\over {x_n(x_nx_{n-1}-c)}},$$  
with an invariant of the form:
$$I(x_n,x_{n-1})=(1+x_{n-1})(1+x_n)(1+{1\over{x_nx_{n-1}-c}}).$$
This equation  was considered by C. Schinas and I. Schinas, and the stability 
was established by using KAM theory, see [L].

Equilibrium $p$ of this equation satisfies $p^4-2cp^2-p+c^2-c=0$, or 
$$(p^2-c)^2=p+c.$$
It is not hard to see that this equation has a positive root $p_+$, which will 
be denoted as $p$ in what follows. The above equilibrium equation implies that 
$p^2>c$. We will be interested in stability of this 
root. Now, we try to apply Theorem A. The necessary conditions for the extremum 
give:
$${\partial I \over \partial x}\equiv (1+y)(1+{1\over{xy-c}}-{{y+xy}\over{(xy-c)^2}})=0,$$
and
$${\partial I \over \partial x}\equiv (1+x)(1+{1\over{xy-c}}-{{x+xy}\over{(xy-c)^2}})=0,$$
which leads to $x=y$, and to the above equilibrium equation. Thus, $x=y=p$. 
This shows that the critical points are exactly the equilibrium points. Now, 
we will check the Hessian at the positive critical point $p$ :

$$H=\left(\matrix{A&B\cr
                  B&A\cr}\right),$$
where 
$$A={\partial^2 I\over \partial x^2}(p,p)={{2p(1+p)}\over {p^2-c}} >0, $$
and

$$B={\partial^2 I\over {\partial x \partial y}}(p,p)=
{{(p^2-c)(p+c)+(c+1)p^2+4cp+c-c^2}\over{(p^2-c)(p+c)}}.$$
Now,  
$$det\; H=A^2 - B^2 = {{(p+1)(p^2-c)+2c^2}\over {(p^2-c)(p+c)}}>0,$$ 
which shows that the invariant $I$ attains its minimum at $(p,p)$. Thus,   
$$\min \{I(x,y): (x,y)\in D\}=I(p,p),$$ 
and so by our theorem $V(x,y)=I(x,y)-I(p,p)$, which shows that $p$ is stable. 
\bigskip
{\bf Example 5.} 
\medskip
So far all our examples have been limited to the equations and related invariants 
to be the rational functions, which is also the case with all examples in the 
literature. Now, we will make the advantage of our results which also allow 
the other types of the functions. In this regard, we consider the special case of 
equation $(3)$ where we take $f_3(y)\equiv 0$ to get 
$$x_{n+1}={{f_1(x_n)}\over {f_2(x_n)}} - g(x_{n-1}). $$
Assuming for simplicity that $f_1(y)$, and $f_2(y)$ are polynomials of the above  
type and taking 
that $a_1=b_1=c_1=e_1=f_1=0$,  and droping the indices in other coefficients  we 
get the following equation 
$$x_{n+1}={{bx_n^2+ex_n+f}\over {ax_n^2+bx_n+c}} - g(x_{n-1}),$$
with the invariant in a special symmetric case:
$$I(x_n,x_{n-1})= ag(x_{n-1})^2 x_n^2 + b_2 (x_n^2 g(x_{n-1}) + x_n g(x_{n-1})^2) + 
c (x_n^2+g(x_{n-1})^2) +$$
$$e g(x_{n-1}) x_n  + f (x_n + g(x_{n-1})).$$
\bigskip
{\bf REMARK 4.} These examples emphasize the  importance of the 
invariants and of finding some general methods for their construction. 
Especially, we are interested in the automatic construction of the invariants and 
corresponding Liapunov functions (if they exist) that can be implemented by 
using some of CAS (Computer Algebra System). In this regard, we have developped 
{\it Mathematica} package, see [MK] which was succesfully tested on all applications 
of this paper. More applications of this package will be presented in the 
coming communications. 

\vfill\eject

\vskip 1cm
\centerline{\r References}
\vskip 1cm
\item{[BB]} {{\bf J.M. Borwein and P.B. Borwein,} {\it Pi and the AGM,}
Wiley, 1987.}
\bigskip
\item{[FJL]} {{\bf J.Feuer, E.Janowski and G.Ladas,} {\it Invariants for Some 
Rational Recursive Sequences with Periodic Coefficients,}  J. Difference Equations 
and Appl. 2(1996), 167-174.}
\bigskip
\item{[GKL]} {{\bf E.A.Grove, V.L.Kocic and G.Ladas,} {\it Classification of 
Invariants for Certain Difference Equations,} Proc. 2nd Int. Conf. on Difference 
Eq. and Appl.,7-11,?}
\bigskip
\item{[GM]} {{\bf I.Gumovski and C.Mira,} {\it Recurences and Discrete Dynamic  
Systems,} Lecture Notes in Mathematics, Springer, 1980.}
\bigskip
\item{[HBQC]} {{\bf F.A.Haggar, G.B. Byrnes, G.R.W. Quispel and H.W.Capel,}
{\it k-Integrals and k-Lie Symmetries in Discrete Dynamical Systems,} (to appear)}
\bigskip
\item{[KP]} {{\bf W.G.Kelley and A.C.Peterson,} {\it Difference Equations,} 
Academic Press, 1991}
\bigskip
\item{[KL]}  {{\bf V.L.Kocic and G.Ladas,} {\it Global Behavior of Nonlinear 
Difference Equations of Higher Order with Applications,} Kluwer Academic 
Publishers, 1993.}
\bigskip
\item{[K]}  {{\bf M.R.S.Kulenovi\'c,} {\it Invariants and Related Liapunov 
Functions for Difference Equations,} Appl. Math. Letters (to appear)}
\bigskip
\item{[L]}  {{\bf G.Ladas,} {\it personal communication.}
\bigskip
\item{[La]}  {{\bf H.Lauwerier,} {\it Fractals,} Princeton University Press, 1991.}
\bigskip
\item{[MK]} {{\bf O.Merino and M.R.S.Kulenovi\'c,} {\it IFL-Invariants and Liapunov  
Functions for Difference Equations in Mathematica,} 1996.}
\bigskip
\item{[M]}  {{\bf C.Mira,} {\it Chaotic Dynamics,} World Scientific, 1987.}
\bigskip
\item{[QRT]} {{\bf G.R.W.Quispel, J.A.G.Roberts and C.J.Thompson,} {\it Integrable 
Mappings and Soliton Equations I,II,} Phys. Lett. A 126 (1988) 419-421; Physica D 
34 (1989) 183-192.}
\item{[RQ]} {{\bf J.A.G.Roberts and G.R.W.Quispel,} {\it Chaos and Time-reversal 
Symmetry. Order and Chaos in Reversible Dynamical Systems} Phys. Rep. 216(1992), 
53-117.}
\bigskip




\end

